\lstset{style=sql, basicstyle=\ttfamily\scriptsize}

\title[Datenbanken] {Datenbanken}

\subtitle{\emoji{floppy-disk} \textbf{\gls{SQL}: Unter-Abfragen}}

\author[Blum, Cyril] % (optional)
{Naoki Peter\inst{1} \and Cyril Blum \inst{2}}

\institute[KZN] % (optional)
{
  \inst{1}%
  Fachschaft Informatik\\
  Kantonsschule Zürich-Nord
  \and
  \inst{2}%
  Fachschaft Informatik\\
  Kantonsschule im Lee
}

\date[\today] % (optional)

\logo{\includegraphics[width=.8cm]{Logos/FreeFlower.pdf}}

\titlegraphic{\includegraphics[width=\textwidth]{"Grundlagen_Info/06_Datenbanken/Figures/title-emojis"}}

\frame{\titlepage}


\begin{frame}[fragile]{Letztes Mal: Murder Mystery}{Abfrage 1}
\begin{lstlisting}
SELECT * FROM crime_scene_report
WHERE city = 'SQL City'
AND date = '20180115'
and type = 'murder'
\end{lstlisting}
\Wider{\begin{figure}[H]
\centering
\includegraphics[width=\linewidth]{"Grundlagen_Info/06_Datenbanken/Figures/MM_hint1"}
\end{figure}}
\end{frame}

\begin{frame}[fragile]{Letztes Mal: Murder Mystery}{Abfrage(n) 2}
\begin{lstlisting}
SELECT *
FROM person JOIN interview
ON interview.person_id = person.id
WHERE address_street_name LIKE 'Franklin%'
AND name like 'Annabel%'
\end{lstlisting}

\begin{lstlisting}
SELECT *
FROM person JOIN interview
ON interview.person_id = person.id
WHERE address_street_name LIKE 'Northwestern Dr'
ORDER BY address_number DESC
LIMIT 1
\end{lstlisting}

\begin{center}
\textbf{Geht es auch in einer Abfrage? $\rightarrow$ Moodle-Aufgaben!}
\end{center}
\end{frame}
\begin{frame}[fragile]{Unterabfragen}{Eine Abfrage in einer Abfrage}
  Die innere Abfrage (in Klammern) wird zuerst ausgeführt, die äussere verwendet ihr Resultat:
\begin{lstlisting}
SELECT name, centimeters
FROM users
WHERE centimeters > (SELECT AVG(centimeters) FROM users)
\end{lstlisting}
  Liefert die innere Abfrage mehrere Werte, verwendet man \lstinline|IN|:
\begin{lstlisting}
SELECT user_id, body
FROM comments
WHERE user_id IN (SELECT id FROM users WHERE city = 'Hamburg')
\end{lstlisting}
\end{frame}

\begin{frame}{Auftrag}{Unterabfragen}
  \begin{itemize}
    \item \readbox{Skript, Abschnitt 2.14}
    \item \aufgabebox{Aufgaben 2.33, 2.34 im Skript}
    \item \aufgabebox{Moodle-Aufgaben zu Unterabfragen}
  \end{itemize}
\end{frame}
